\section{Binary-fibre transversals and equal-projection obstruction}
\label{sec:binary-transversals}

Let $Q:A\times C\to S$ be any Latin square of order $q$, and let
$h:A\times C\to\mathbb F_2$ be arbitrary. The table
\[
 B((a,x),(b,y))=(Q(a,b),x+y+h(a,b))
\]
is Latin of order $2q$. All bit arithmetic is in $\mathbb F_2$.
These are fibres of size two over an order-$q$
quotient, not an assumed quotient onto an order-two square.
A \emph{simple two-plex} $P$ is a set of $2q$ distinct quotient
cells, meeting each quotient row, column and symbol twice.

Every physical transversal of $B$ projects to such a $P$.
Indeed, a repeated quotient cell would use opposite row bits and
opposite column bits, hence equal symbol bits, repeating a physical
symbol. On the cells of $P$ form three perfect matchings by equal
quotient row, column and symbol. They are edge-disjoint by Latinness.
Let $G_P$ be their union and let $c(P)$ be its number of components.
For $e=(a,b,Q(a,b))$ write $h_e=h(a,b)$.
Each matching restricts to a perfect matching on every component,
whose number of vertices is therefore even.

\begin{theorem}[Component-parity lift criterion and exact count]
\label{thm:binary-lift-count}
The simple two-plex $P$ lifts to a physical transversal of $B$ if and
only if, for every component $K$ of $G_P$,
\[
 \sum_{e\in V(K)}h_e=|V(K)|/2\pmod2.
\]
If consistent, it has exactly $2^{q+c(P)}$ distinct labelled lifts.
\end{theorem}
\begin{proof}
For each selected cell $e=(a,b,Q(a,b))$, choose row and column bits
$x_e,y_e$, and set $z_e=x_e+y_e+h_e$. The transversal equations are
\[
\begin{aligned}
 x_e+x_f&=1 &&\text{on a row edge},\\
 y_e+y_f&=1 &&\text{on a column edge},\\
 x_e+x_f+y_e+y_f&=1+h_e+h_f &&\text{on a symbol edge}.
\end{aligned}
\]
Summing within a component cancels all variables and gives exactly
the displayed parity condition.

There are $4q$ variables and $3q$ equations. A linear dependency with
coefficients $u_a,v_b,w_s$ requires $u_a=w_s=v_b$ at each selected
cell. Equality propagates over the three matchings, so a dependency
is precisely a choice of one constant per connected component.
The left kernel therefore has dimension $c(P)$ and the matrix has
rank $3q-c(P)$. The component conditions test a basis of that left
kernel, so they are sufficient as well as necessary. The solution
space has dimension $q+c(P)$. Distinct solutions give distinct
physical transversals, and every transversal supplies one solution.
Independent bit relabelling of any quotient row, column or symbol
adds twice to each component sum, so the criterion is label-invariant.
\end{proof}

Standard double prolongation along two disjoint transversals $T_0,T_1$
replaces each selected triple $(r,c,s)\in T_t$ by
$(r,c,N_t),(r,N_t,s),(N_t,c,s)$ and fills the new $2\times2$ corner
with either Latin order-two table.

\begin{theorem}[Equal-projection obstruction in three views]
\label{thm:same-projection-obstruction}
For every Latin $Q$ of order $q\geq2$ and every $h$, a standard double
prolongation of $B$ along disjoint transversals with the same simple
two-plex projection is not FFF, for either corner.
No FFF hypothesis on $Q$ or $B$ is needed.
\end{theorem}
\begin{proof}
At each projected cell compare the two lifts and put
$d_x=x'_e+x_e$, $d_y=y'_e+y_e$, $d_z=z'_e+z_e=d_x+d_y$.
The transversal equations make $d_x$ constant on row-matching edges,
$d_y$ on column edges, and $d_z$ on symbol edges.
Disjointness forbids $d_x=d_y=0$.

Suppose $d_x=0$ on one quotient row. Write its physical rows as
$r,r^*$ and let $\beta$ interchange the two columns of every column
fibre. If $T_0$ selects columns $u,v$ in $r,r^*$, then $T_1$ selects
$\beta(u),\beta(v)$ there. The two $\beta$-orbits are distinct.
Originally the induced row permutation is $\beta$. After prolongation
its arrows contain
\[
 u\longmapsto v\longmapsto N_1\longmapsto u,\qquad
 \beta(u)\longmapsto\beta(v)\longmapsto N_0
       \longmapsto\beta(u).
\]
These follow directly by matching the replaced symbols and the displaced
old entries; other column fibres stay transpositions. The exact pair
type is $3^2 2^{q-2}$, independently of the corner. Thus row-F requires
$d_x=1$ on every selected cell.

The triple description of prolongation commutes with permutations
of row, column and symbol coordinates. Each order-two corner has
equation $i+j+k=\kappa$ over $\mathbb F_2$ and is likewise symmetric.
The binary-fibre equation remains of the same form in every view,
with the transported quotient and the same projected triple set.
The identical paired-line argument therefore makes column-F require
$d_y=1$ everywhere and symbol-F require $d_z=1$ everywhere.
But $d_z=d_x+d_y=0$ if the first two requirements hold. All three
views cannot be F.
\end{proof}

\begin{proposition}[Canonical first transversal, arbitrary second]
\label{prop:free-first-orbits}
The four global bit autotopies, with $u,v\in\mathbb F_2$,
$(a,x)\mapsto(a,x+u)$, $(b,y)\mapsto(b,y+v)$,
$(s,z)\mapsto(s,z+u+v)$ act freely on physical transversals.
With numeric labels $2a+x,2b+y$, each orbit has exactly one column
tuple $f$ with $\lfloor f(0)/2\rfloor<\lfloor f(1)/2\rfloor$ and
$f(0)$ even. Restricting only the first transversal to this form
loses no standard double-prolongation FFF candidate of the fixed base.
\end{proposition}
\begin{proof}
The transformed tuple is $f'(r)=f(r\mathbin{\mathrm{xor}}u)
\mathbin{\mathrm{xor}}v$. A nonzero $v$ with $u=0$ changes every
column. A fixed tuple with $u=1$ would put both physical rows of
each quotient row in the same quotient column, which a transversal
cannot do. Thus the action is free. The two quotient columns at
rows $0,1$ are distinct: a unique $u$ orders them and a unique $v$
makes the first column even.
Extend the three coordinate bijections by fixing $N_0,N_1$.
They transport every replacement triple and fix either corner.
Once the first is normalized, they bijectively transport all disjoint
second transversals. The second must remain unrestricted; demanding
that it too be canonical would need another justification.
\end{proof}

\paragraph{A count is not mate coverage.}
For the frozen non-group FFF16 binary-fibre base \texttt{source21 comm}
in \nolinkurl{evidence/data/source21_base.json}, the complete quotient
census is $181768$ simple two-plexes, of which $138486$ lift.
The lift formula gives $72474624$ labelled first transversals and
$18118656$ representatives under the four autotopies, not main classes.
An independently implemented count agrees. At this edition's cutoff,
only $18$ selected supports, comprising $9216$ labelled firsts, have
separate complete arbitrary-mate exclusions: two previously closed
supports and sixteen newly checked whole supports.
The remaining $138468$
liftable support domains are not excluded by that package.
Different-projection pairs lie outside
Theorem~\ref{thm:same-projection-obstruction}. That theorem applies to
every binary-fibre base satisfying its hypotheses; the finite census
and arbitrary-mate exclusions above concern only this specified base.
Unrestricted FFF18 remains open.
